\documentclass[11pt]{article}

\usepackage[utf8]{inputenc}
\usepackage[T1]{fontenc}
\usepackage[margin=0.9in]{geometry}
\usepackage{amsmath,amssymb,amsfonts}
\usepackage{mathtools}
\usepackage{bm}
\usepackage{booktabs}
\usepackage{array}
\usepackage[hidelinks]{hyperref}
\usepackage{enumitem}
\setlist{nosep,leftmargin=1.4em}

\newcommand{\clip}{\operatorname{clip}_{[0,1]}}
\newcommand{\Real}{\mathbb{R}}
\DeclareMathOperator*{\argmin}{arg\,min}
\newcommand{\half}{\tfrac{1}{2}}

\title{\textbf{Participation Diffusion Dynamics in\\ Influence-Driven Networks}}
\author{Hudson Dong \thanks{Class of 2027, Stuyvesant High School, 345 Chambers St, New York, NY 10282. Email: \texttt{hdong70@stuy.edu}}}
\date{}

\begin{document}
\maketitle

\begin{abstract}
\noindent
Student-club participation results from socially embedded, conditionally renewable action: a member may attend one week, skip the next, and perhaps return later. Each act of participation reshapes the social environment that governs future attendance and intentions. Instead of using the one-shot ``adopt or not'' event that classical diffusion models assume, I develop an agent-based, discrete-time model of participation on a weighted, directed social network in which (i) participation is a probabilistic, repeated behavior generated by latent internal states: motivation, loyalty, fatigue, and three slower modulating states; (ii) the effective decay of motivation is endogenous, slowed by loyalty and accelerated by fatigue; and (iii) social influence itself co-evolves with activity, with visibility and persuasiveness treated as distinct, dynamic factors. The model nests the Linear Threshold and Independent Cascade processes as limiting cases. I then treat the system as a controllable dynamical system: intervention levers (events, visibility amplification, capacity, early scaffolding) enter only through existing channels, which enables reverse planning of minimum-cost interventions, an explicit feasibility frontier that separates structural infeasibility from optimization failure, closed-loop receding-horizon control, and responsible-design constraints on load, equity, and robustness. The framework converts a qualitative contrast, such as the stabilization of one club against the collapse of a comparable group, into parameter regimes, intervention levers, and costed plans. An accompanying browser simulator implements the same equations as a computational readout: users may inspect a single randomized sample trajectory, average many randomized network realizations, or reverse engineer plans using the underlying behavioral model.
\end{abstract}

\section{Introduction}

\subsection{Background and Motivation}
Student clubs are a pervasive organizational form in secondary education. They are voluntary, peer-mediated, and sustained through regular commitment rather than formal incentives, which places them in an analytically awkward middle ground: participation is neither a one-shot adoption decision nor a mandated, continuously enforced behavior like class attendance. It unfolds instead as repeated effort under time scarcity, social visibility, and heterogeneous skill readiness. These features make participation a dynamic, state-dependent, and socially embedded process.

The model is motivated by a sharply controlled contrast. Specifically, within one high school over comparable periods and under identical administrative constraints, the same organizer launched two clubs using similar recruitment channels. A fitness-oriented club diffused rapidly and stabilized into persistent weekly participation; a composition-oriented club attracted early attention but decayed and eventually dissolved. The divergence is hard to attribute to leadership, network size, or institutional support: the two clubs were embedded in overlapping social networks and exposed to similar informational environments. The contrast instead points to \emph{endogenous} participation dynamics---early experiences appear to have reshaped motivation differently across domains, and peer signals translated into reinforcement in one case but not the other. Comparable patterns appear in online and offline settings, where large early inflows do not guarantee persistence and where reinforcement, rather than reach, predicts long-term survival (Centola, 2010; Kairam et al., 2012).

This observation exposes a measurement problem that runs through diffusion research: observed diffusion curves conflate at least two mechanisms. One is social transmission across a network; the other is within-person persistence, shaped by perceived progress, fatigue, and repeated exposure. Field and experimental evidence shows that social influence can be substantial but is heterogeneous and reinforcement-dependent rather than purely exposure-driven (Aral \& Walker, 2012; Bond et al., 2012). Distinguishing the two mechanisms is the central modeling goal of this paper.

\subsection{Limits of Standard Diffusion Models}
Classical diffusion models on networks treat adoption as a binary event on a fixed graph. Threshold and stochastic-cascade models remain foundational and provide the algorithmic basis for influence maximization (Watts, 2002; Kempe et al., 2003), and they have been extended to weighted, heterogeneous, and temporal networks. Yet they abstract away from a defining feature of clubs: participation is \emph{renewable}, and it reshapes the system that governs future participation. Visibility can rise after posting; persuasiveness can change after repeated co-activity; bursty interaction patterns can slow or accelerate spreading even when the static structure looks favorable (Holme \& Saram\"aki, 2012).

A second limitation concerns reinforcement. For behaviors that require effort, skill, or public exposure, a single exposure is often insufficient: social reinforcement from multiple peers can dominate, and clustered ties can become an advantage rather than a bottleneck (Centola, 2010; Ugander et al., 2012). This challenges models that treat influence as additive and memoryless, shifting what ``good network structure'' means for behavioral diffusion. Both limitations bear directly on the motivating contrast: the composition club plausibly operated in a regime where early participation did not reliably generate reinforcing experience, whereas fitness engagement offered lower entry barriers and immediate feedback. The implication is structural: diffusion outcomes depend not only on who is connected to whom, but on how internal participation states evolve over time.

\subsection{Research Questions}
The paper develops a controllable participation-diffusion framework organized around three questions. \emph{First}, under what conditions does participation become self-sustaining rather than transient? This concerns the dynamic balance between reinforcement and decay and whether repeated involvement builds stabilizing internal stocks that protect against dropout. \emph{Second}, how do domain-specific participation demands interact with network influence to produce divergent trajectories under similar external conditions? This treats ``fitness versus composition'' as a difference in constraints, feedback delay, and visibility rather than in leadership or recruitment effort. \emph{Third}, how can interventions transform a failing trajectory into a robust one at minimal cost? This motivates an inverse problem: given a target participation trajectory, recover a feasible intervention path subject to resource and feasibility constraints.

\subsection{Contributions and Scope}
Section~\ref{sec:model} introduces an agent-based dynamical system on a weighted, directed network. Individuals carry latent internal states governing repeated participation; motivation evolves with an effective decay endogenously shaped by loyalty and fatigue. Furthermore, external events, internal reinforcement from perceived progress, and social influence jointly determine transitions. Thus, observable actions emerge probabilistically rather than through deterministic thresholds. Social influence is itself endogenous, with visibility and persuasiveness distinguished and influence weights evolving through co-activity. Section~\ref{sec:control} treats the model as a controllable diffusion system: intervention levers enter the dynamics explicitly, enabling reverse planning that infers cost-minimizing schedules subject to capacity and feasibility constraints. Together, these extensions apply control-of-spreading ideas to settings where persistence---not adoption count---is the objective (Nowzari et al., 2016; Preciado et al., 2014). Section~\ref{sec:feasibility} delineates the limits of the approach, clarifying which failures reflect poor design and which arise from structural infeasibility.

The contribution is twofold. Substantively, the model provides a unified mechanism for participation diffusion that integrates internal-state persistence, social reinforcement, and network co-evolution. Methodologically, it offers a control-oriented framework that converts qualitative contrasts such as ``one club succeeds while another fails'' into testable parameter regimes, intervention levers, and costed implementation paths.

\subsection{Roadmap}
The remainder of the paper is organized as follows. Section~\ref{sec:related} situates the model against threshold and cascade diffusion, behavioral and habit-formation theory, empirical work on group survival, co-evolutionary networks, and the control-of-spreading literature. Section~\ref{sec:model} develops the model in full: the probabilistic action map, the motivation dynamics with endogenous decay, the loyalty and fatigue stocks, external and internal reinforcement, the co-evolving influence structure, capacity, attrition, and the modulating states. It then gives a mean-field fixed-point analysis showing how bistability, and hence the collapse-versus-stabilization dichotomy, arises directly from the equations. Section~\ref{sec:control} recasts the system as a controllable diffusion process, formulates the cost-minimizing inverse problem, specifies two domain regimes and three reference scenarios, and traces each scenario at the level of the underlying mechanism. Section~\ref{sec:feasibility} establishes the feasibility frontier separating structural infeasibility from optimization failure, develops closed-loop receding-horizon control with state estimation, adds responsible-design constraints on individual load, subgroup equity, and robustness, and describes the accompanying simulator as a readout layer. Section~\ref{sec:discussion} discusses assumptions and limitations, and Section~\ref{sec:future} outlines an empirical calibration and validation program.

\section{Related work}\label{sec:related}

\paragraph{Classical Diffusion on Networks.}
Diffusion in networks has a long tradition across sociology, physics, and computer science. The threshold model, first introduced by Granovetter then formalized in network settings by Watts, models adoption as a binary transition once peer influence exceeds an individual threshold (Granovetter, 1978; Watts, 2002), while the independent-cascade family treats adoption as a stochastic process in which active neighbors independently trigger activation (Kempe et al., 2003). A closely related physics tradition studies spreading through epidemic-style dynamics, where degree heterogeneity can qualitatively alter thresholds and persistence (Pastor-Satorras \& Vespignani, 2001; Pastor-Satorras et al., 2015). Despite extensions to heterogeneous thresholds, weighted edges, and temporal networks, the core abstraction is unchanged: adoption is a one-time event, and once adopted an individual remains active. Effort costs, learning curves, and fatigue are not represented, so these models suit information spread better than activities requiring repeated voluntary engagement. The distinction is conceptual: classical diffusion concerns \emph{states}; participation concerns \emph{behaviors over time}.

\paragraph{Behavioral and Reinforcement-based Perspectives.}
A separate literature emphasizes the internal processes governing persistence. Self-determination theory distinguishes intrinsic from extrinsic regulation and highlights internalization in sustained behavior (Ryan \& Deci, 2000), while habit-formation research shows that repeated action under stable conditions reduces cognitive effort and increases persistence, but slowly and with sensitivity to disruption (Lally et al., 2010). These accounts explain why individuals disengage after initial adoption and why early experiences are decisive, but they are typically individual-level: social influence enters implicitly and network structure is rarely modeled. They therefore do not explain why identical motivational profiles yield different outcomes under different social embeddings, nor how individual reinforcement aggregates into collective trajectories.

\paragraph{Empirical Evidence on Persistence and Survival.}
Studies of online communities provide direct evidence that early participation dynamics matter more than reach: across thousands of groups, growth rate and early retention strongly predict survival, and many groups with large early inflows collapse from weak internal engagement (Kairam et al., 2012). Experimental work shows that behaviors requiring reinforcement spread differently from simple contagions; adoption depends on exposure from multiple peers, so clustered structures can outperform random ones for behavioral diffusion even as they slow simple cascades (Centola, 2010). Diffusion success thus depends on whether participation generates reinforcing feedback and whether social exposure amplifies or dampens it, not on size or connectivity alone.

\paragraph{Endogenous Influence and Co-evolution.}
Social influence is not static: tie strength, credibility, and visibility evolve through interaction, and repeated co-activity strengthens ties and reshapes future pathways. This connects to adaptive and co-evolutionary networks, in which node states and topology evolve jointly (Gross \& Blasius, 2008), and where coupling state dynamics to edge dynamics can change equilibrium structure and stability---as in nonequilibrium phase behavior when opinions co-evolve with rewiring (Holme \& Newman, 2006). Most diffusion models nonetheless treat influence weights as fixed, and when adaptation is modeled it is usually decoupled from participation persistence---limiting their ability to explain why some groups thicken internally while others fragment.

\paragraph{From Diffusion to Control.}
Influence maximization introduces optimization into diffusion, but typically optimizes over initial seeds under fixed rules, and over adoption counts rather than sustained participation (Kempe et al., 2003). An adjacent line treats spreading as a dynamical process to be shaped through control and resource allocation: surveys formalize the analysis and controllability of spreading on networks (Nowzari et al., 2016), and optimization formulations allocate preventive and corrective resources under budget constraints, making feasibility and cost explicit (Preciado et al., 2014). Three gaps remain: transient engagement is rarely distinguished from stable participation; intervention design is largely heuristic and seldom embedded in the diffusion equations as explicit control inputs; and feasibility limits are not addressed, so failure cannot be attributed to design versus structural infeasibility. These gaps are acute in youth settings, where participation is voluntary, resources are constrained, and ethics limit intervention intensity.

\paragraph{Positioning.}
The present study integrates motivational reinforcement, fatigue, and loyalty into a network-diffusion system with endogenous, co-evolving influence and probabilistic, repeated participation, and treats diffusion as a controllable dynamical system. This enables not only explanation and prediction but inverse planning under cost and feasibility constraints, connecting diffusion theory to organizational design.

% =====================================================================
\section{The model}\label{sec:model}

\begin{table}[h]
\centering\footnotesize
\setlength{\tabcolsep}{4pt}
\renewcommand{\arraystretch}{1.2}
\begin{tabular}{@{}ll@{\hspace{1.2em}}ll@{}}
\toprule
\multicolumn{4}{@{}l}{\textit{Per-node states (time-varying)}}\\
$m_i(t)$ & motivation & $S_i(t)$ & social saturation\\
$L_i(t)$ & loyalty (habit/commitment) & $\Phi_i(t)$ & stability weight (volatility damping)\\
$F_i(t)$ & fatigue (short-term depletion) & $\mathrm{act}_i(t)\!\in\!\{0,1\}$ & active flag\\
$E_i(t)$ & emotional state (affect) & $\mathrm{cred}_i(t)$ & credibility (derived)\\
\midrule
\multicolumn{4}{@{}l}{\textit{Per-node actions}}\\
$a_i(t)$ & attendance (Bernoulli/prob.) & $p_i(t)$ & posting / sharing\\
$\tilde a_i(t)$ & effective attendance (post-capacity) & $y_i(t)\!\in\![0,T_i]$ & training volume\\
\midrule
\multicolumn{4}{@{}l}{\textit{Per-node parameters}}\\
$\alpha_{0,i}\!\in\!(0,1)$ & baseline self-decay & $w^{P}_i\!>\!0$ & peer susceptibility\\
$b_i\!\in\![0,1]$ & spreading propensity & $\theta^a_i,\theta^p_i$ & action thresholds\\
$\kappa^a_i,\kappa^p_i\!>\!0$ & logistic slopes & $D_i\!\in\![0,1]$ & domain preparedness\\
$V^{ex}_i,V^{ev}_i\!\in\![0,1]$ & broadcast / event visibility & $c_i(t)\!\in\![0,1]$ & availability\\
$T_i(t),\bar y_i$ & volume cap, nominal volume & $s^{\mathrm{role}}_i(t)$ & role salience\\
$\beta_i\!>\!0$ & internal-reinf.\ cap & $\delta_i\!>\!0$ & internal-reinf.\ sensitivity\\
$\chi_i,\chi^{\mathrm{viab}}_i\!\ge\!0$ & skip / below-quorum penalty & $\rho_i\!>\!0$ & soft-saturation (resistance)\\
$h_{0,i}\!\ge\!0$ & baseline dropout hazard & $\gamma_1,\gamma_2,\gamma_3\!>\!0$ & hazard elasticities\\
$\xi_1,\xi_2\!>\!0$ & reactivation gains & $\zeta^0_i,\zeta^1_i,\zeta^2_i$ & persuasiveness coeffs\\
\midrule
\multicolumn{4}{@{}l}{\textit{Global / coupling}}\\
$\eta_L,\eta_F\!\ge\!0$ & decay sensitivities & $\lambda_L,\lambda_F\!\in\!(0,1]$ & loyalty / fatigue rates\\
$\varphi_1,\varphi_2,\varphi_3$ & loyalty channel weights & $\psi_1,\psi_2,\psi_3$ & fatigue channel weights\\
$\mu\!\in\!(0,1],\eta_1,\eta_2$ & Hebbian rate, gains & $\mu_{\mathrm{mix}}\!\ge\!0$ & pull toward channel mix\\
$\pi_o,\pi_f\ (\pi_o\!+\!\pi_f\!=\!1)$ & online / offline mix & $w^c_1,w^c_2,\Delta$ & credibility weights, window\\
$\kappa_g\!>\!0$ & transmission sharpness & $\beta_E\!\ge\!0$ & emotion\,$\to$\,action gain\\
$r^+(t)\!\in\![0,1]$ & positive-narrative signal & $K(t)\!>\!0$ & capacity\\
$N_{\min}$ & minimum viable participation & $\{(\tau_k,A_k)\},\lambda_R$ & events, decay rate\\
$\mathrm{sim}_{ij}\!\in\![0,1]$ & similarity & $\epsilon_i(t)$ & zero-mean shock\\
\bottomrule
\end{tabular}
\caption{Notation. Coupling gains $\beta_E,\kappa_{EF},\kappa_{SL},\kappa_\Phi$ for the modulating states (Sec.~\ref{sec:modulating}) default to zero, recovering the core model.}
\label{tab:notation}
\end{table}

\subsection{Setup and notation}\label{sec:setup}
The system evolves on a weighted, directed graph $G=(V,E,W_t)$ with $|V|=n$ students and discrete weeks $t=0,1,\dots,T$. The directed edge $j\to i$ carries weight $\omega_{ij}(t)\in[0,1]$, the influence of sender $j$ on receiver $i$. Every node holds latent states; observable actions are random draws from those states. Unless noted, all quantities lie in $[0,1]$; $\clip$ projects onto $[0,1]$; $[x]_+=\max(x,0)$ and $[x]_-=\max(-x,0)$; $\sigma(x)=1/(1+e^{-x})$ is the logistic link; $\mathbf 1\{\cdot\}$ is the indicator; and $\epsilon_0$ is a small constant guarding divisions. Table~\ref{tab:notation} collects the symbols. Two aggregates and the emotion-shifted motivation recur throughout:
\begin{equation*}
\bar a_{N(i)}(t)=\frac{\sum_j \omega_{ij}(t)\,\tilde a_j(t)}{\sum_j \omega_{ij}(t)+\epsilon_0},
\qquad
\hat m_i(t)=m_i(t)+\beta_E\big(E_i(t)-\half\big).
\end{equation*}

\subsection{Observable actions and the S-shaped response}\label{sec:actions}
Attendance $a_i$ and posting $p_i$ are probabilistic realizations of emotion-shifted motivation through a logistic map:
\begin{equation}\tag{A}\label{eq:actions}
P^a_i(t)=\sigma\!\big(\kappa^a_i(\hat m_i(t)-\theta^a_i)\big),\qquad
P^p_i(t)=\sigma\!\big(\kappa^p_i(\hat m_i(t)-\theta^p_i)\big).
\end{equation}
In \emph{stochastic} mode, $a_i(t)\sim\mathrm{Bernoulli}(P^a_i)\cdot\mathrm{act}_i$ and $p_i(t)\sim\mathrm{Bernoulli}(P^p_i)\cdot\mathrm{act}_i$; in \emph{deterministic} mode, $a_i(t)=P^a_i(t)\cdot\mathrm{act}_i$. As $\kappa\to\infty$ the response approaches a hard threshold $a_i=\mathbf 1\{\hat m_i\ge\theta^a_i\}$; an optional hysteresis variant activates at $\theta^{\uparrow}$ and holds until $\hat m_i<\theta^{\downarrow}<\theta^{\uparrow}$, producing the commitment ``stickiness'' typical of established members. Separate thresholds and slopes for attending versus posting let the visible social-media footprint grow faster or slower than physical participation---an empirically important asymmetry, since posting is cheap relative to attendance.

\subsection{Motivation and self-decay}\label{sec:motivation}
Raw peer input is the weighted sum of neighbors' transmitted motivation, then soft-saturated to remain bounded and differentiable:
\begin{equation}\tag{15}\label{eq:softsat}
X_i(t)=\sum_j \omega_{ij}(t)\,b_j\,g\!\big(m_j(t)\big),\quad g(x)=\tanh(\kappa_g x);
\qquad \tilde X_i(t)=\frac{X_i(t)}{1+\rho_i X_i(t)}.
\end{equation}
Motivation then evolves as retention plus a fluctuating drive comprising peer influence, external and internal reinforcement, and noise:
\begin{equation}\tag{1}\label{eq:mot}
m_i(t{+}1)=\clip\!\Big[(1-\alpha_i(t))\,m_i(t)+G(\Phi_i(t))\cdot\big(w^{P}_i(1-S_i(t))\tilde X_i(t)+c_i(t)R_i(t)+B_i(t)+\epsilon_i(t)\big)\Big],
\end{equation}
where the stability gain damps the drive when the trajectory has been calm ($\Phi\to1$) and amplifies it when volatile ($\Phi\to0$):
\begin{equation*}
G(\Phi)=1+\kappa_\Phi(1-2\Phi),\qquad \kappa_\Phi\in[0,1).
\end{equation*}
Setting all $\omega_{ij}=0$ collapses \eqref{eq:mot} to independent exponential decay $m_i(t{+}1)=(1-\alpha_i)m_i$, so every isolated student eventually disengages. With moderately active peers the peer sum offsets self-decay and motivation stabilizes at an intermediate level (a club with steady but modest activity); when an event raises $R_i$, the S-shaped response makes many attendance probabilities jump within one or two weeks, creating a collective surge whose persistence after the pulse depends on whether loyalty has accumulated enough to retain the new participants. These three regimes---independent decay, peer-maintained plateau, and pulse-driven surge---are the qualitative vocabulary the rest of the model formalizes.

The effective decay is itself endogenous: loyalty slows it, fatigue accelerates it, negative affect inflates effective fatigue, and saturation weakens loyalty's protection:
\begin{equation}\tag{2}\label{eq:decay}
\begin{aligned}
\alpha_i(t)&=\clip\!\Big[\alpha_{0,i}\big(1-\eta_L(1-\kappa_{SL}S_i(t))L_i(t)\big)\big(1+\eta_F\hat F_i(t)\big)\Big],\\
\hat F_i(t)&=\clip\!\big[F_i(t)+\kappa_{EF}(\half-E_i(t))\big].
\end{aligned}
\end{equation}
The multiplicative form keeps $\alpha_i$ positive and bounded and lets loyalty and fatigue moderate one another rather than merely add. The decay rate is the model's central coupling: it is where the slow stabilizers (loyalty) and destabilizers (fatigue) feed back onto the fast variable (motivation), and it is the variable a successful intervention must ultimately move.

\subsection{Loyalty and fatigue}\label{sec:loyalty}
Both stocks update toward behavior-based targets. Their self-reinforcement channel is \emph{effective attendance} $\tilde a_i$, not the decay rate:
\begin{equation}\tag{3}\label{eq:loyalty}
\begin{aligned}
L_i(t{+}1)=\clip\!\Big[L_i(t)+\lambda_L\big(&\varphi_1\tilde a_i(t)+\varphi_2\,s^{\mathrm{role}}_i(t)(1-\kappa_{SL}S_i(t))\bar a_{N(i)}(t)\\
&+\varphi_3\,\nu_E(E_i)\,r^+(t)-L_i(t)\big)\Big].
\end{aligned}
\end{equation}
\begin{equation}\tag{4}\label{eq:fatigue}
\begin{aligned}
F_i(t{+}1)=\clip\!\Big[F_i(t)+\lambda_F\big(&\nu_E(E_i)\psi_1\,y_i(t)+\nu_E(E_i)\psi_2\,\bar a_{N(i)}(t)\\
&+\psi_3(1-\tilde a_i(t))-F_i(t)\big)\Big].
\end{aligned}
\end{equation}
where the mood multiplier raises fatigue accrual and blunts narrative response under negative affect:
\begin{equation*}
\nu_E(E_i)=1+\beta_{EF}\big(\half-E_i(t)\big).
\end{equation*}
The recovery term $\psi_3(1-\tilde a_i)$ is positive only when resting, producing effort--rest cycles rather than monotonic exhaustion. Loyalty is a stock of positive inertia and fatigue a stock of negative inertia; both feed back into the decay rate \eqref{eq:decay}. Their steady states---sustained attendance and positive narratives raise loyalty, sustained exertion raises fatigue---define each student's personalized ``half-life'' of motivation, and the gap between these half-lives across domains is what the regimes of Section~\ref{sec:control} encode.

\subsection{External and internal reinforcement}\label{sec:reinforcement}
External events enter as exponentially decaying pulses (half-life $\ln 2/\lambda_R$), gated by availability $c_i(t)$ inside \eqref{eq:mot} and scaled by saturation and mood:
\begin{equation}\tag{5}\label{eq:external}
R_i(t)=(1-S_i(t))\,\nu_E(E_i)\,V^{ev}_i\!\!\sum_{k:\,\tau_k\le t}\! A_k\,e^{-\lambda_R(t-\tau_k)}.
\end{equation}
Internal reinforcement converts perceived progress into motivation, with a minimum-viability sign flip: below quorum, reinforcement becomes a flat penalty even for prepared members, capturing the discouragement of an under-attended session:
\begin{equation}\tag{6}\label{eq:internal}
B_i(t)=\nu_E(E_i)\cdot\gamma_B(t)\cdot
\begin{cases}
\beta_i\big(1-e^{-\delta_i y_i(t)}\big)-\chi_i\,\mathbf 1\{y_i(t)=0\}, & A(t)\ge N_{\min},\\[3pt]
-\,\chi^{\mathrm{viab}}_i, & A(t)<N_{\min},
\end{cases}
\end{equation}
with training volume and its half-saturation point
\begin{equation*}
y_i(t)=\min(T_i(t),\bar y_i)\,\tilde a_i(t),\qquad y_i^{1/2}=\ln 2/\delta_i.
\end{equation*}
The concave gain $\beta_i(1-e^{-\delta_i y_i})$ reproduces diminishing returns to effort; the discrete penalty $-\chi_i\mathbf 1\{y_i=0\}$ reproduces streak effects; and the quorum dependence is precisely what makes a composition-type domain (high skill, delayed feedback, quorum-dependent) collapse while a fitness-type domain (low barriers, immediate feedback, productive at small sizes) stabilizes. Here $D_i$ enters through the calibration of $\beta_i,\delta_i,N_{\min}$: domains with low preparedness have small $\beta_i$, large $\delta_i$ (slow progress), and high $N_{\min}$.

\subsection{Social-influence structure}\label{sec:influence}
The influence weight decomposes into a visibility factor times a persuasiveness factor:
\begin{equation}\tag{7}\label{eq:omega}
\omega_{ij}(t)=\clip\!\big[v_{ij}(t)\,s_{ij}(t)\big].
\end{equation}
Persuasiveness combines baseline trust, similarity, and a dynamic credibility that rises with the sender's loyalty and recent attendance:
\begin{align}
s_{ij}(t)&=\zeta^0_i+\zeta^1_i\,\mathrm{sim}_{ij}+\zeta^2_i\,\mathrm{cred}_j(t),\tag{9}\label{eq:persuade}\\
\mathrm{cred}_j(t)&=w^{c}_1 L_j(t)+w^{c}_2\,\tfrac1\Delta\!\!\sum_{\tau=t-\Delta+1}^{t}\!\!\tilde a_j(\tau).\tag{10}\label{eq:cred}
\end{align}
Visibility blends a structural channel mix (offline and online exposure) with a Hebbian co-activity term; this is where the broadcaster's expressive visibility $V^{ex}_j$---the channel the control $\gamma_v$ amplifies---enters:
\begin{equation}\tag{8}\label{eq:vmix}
v^{\mathrm{mix}}_{ij}(t)=\pi_o\,\mathrm{online}_{ij}(t)+\pi_f\,\mathrm{offline}_{ij}(t),\qquad \pi_o+\pi_f=1,
\end{equation}
\begin{equation}\tag{11}\label{eq:vupdate}
v_{ij}(t{+}1)=\clip\!\Big[(1-\mu)v_{ij}(t)+\mu\big(\eta_1\,\tilde a_i(t)\tilde a_j(t)+\eta_2\,V^{ex}_j\,p_j(t)\big)+\mu_{\mathrm{mix}}\big(v^{\mathrm{mix}}_{ij}(t)-v_{ij}(t)\big)\Big].
\end{equation}
Co-attendance thickens ties symmetrically; $j$'s posting thickens the edge $i\!\leftarrow\! j$ in proportion to $j$'s broadcast strength; and the $\mu_{\mathrm{mix}}$ term keeps data-driven exposure relevant. Repeated application drives $W_t$ toward block (community) structure, and the update is contractive for $\mu(\eta_1+\eta_2)<1$. Two limiting cases are useful: the pure-Hebbian model is recovered by setting $\mu_{\mathrm{mix}}=0$ and $s_{ij}\equiv1$, and freezing $v_{ij}$ recovers a static-graph diffusion. Separating visibility from persuasiveness matters because the two respond to different levers: visibility is amplified by broadcasting and co-activity, persuasiveness by accumulated credibility.

\subsection{Capacity and resource constraints}\label{sec:capacity}
When intended attendance exceeds capacity, realized attendance is scaled down proportionally; a smooth, differentiable variant is used inside optimization:
\begin{equation}\tag{12}\label{eq:capacity}
\tilde a_i(t)=a_i(t)\cdot\min\!\Big(1,\ \tfrac{K(t)}{\sum_j a_j(t)+\epsilon_0}\Big)
\qquad\Big[\text{smooth: }\ \tilde a_i(t)=\tfrac{a_i(t)}{1+(\sum_j a_j(t))/K(t)}\Big].
\end{equation}
Effective attendance $\tilde a_i$---not intended $a_i$---is what propagates into \eqref{eq:loyalty}, \eqref{eq:fatigue}, \eqref{eq:internal}, \eqref{eq:cred}, and \eqref{eq:vupdate}, so congestion has downstream psychological and structural consequences and growth is self-limiting. This is why capacity bounds matter even when the objective is growth: clipped participation slows the very accumulation of loyalty and credibility that would otherwise reduce decay and thicken $W_t$.

\subsection{Attrition and reactivation}\label{sec:attrition}
Participation is not an absorbing state. Active members drop out with a hazard that falls with motivation and loyalty and rises with fatigue:
\begin{equation}\tag{13}\label{eq:hazard}
h_i(t)=h_{0,i}\,e^{-\gamma_1 m_i(t)-\gamma_2 L_i(t)+\gamma_3 F_i(t)},\qquad d_i(t)\sim\mathrm{Bernoulli}\!\big(1-e^{-h_i(t)}\big).
\end{equation}
On dropout, $\mathrm{act}_i\!\leftarrow\!0$, $a_i=p_i=0$, and outgoing visibility decays by a factor $\kappa_{\mathrm{drop}}<1$, modeling the social disappearance of inactive members. Inactive members return under event exposure and neighbor activity:
\begin{equation}\tag{14}\label{eq:react}
P^{\mathrm{react}}_i(t)=1-e^{-\xi_1 R_i(t)-\xi_2\bar a_{N(i)}(t)}.
\end{equation}
The multiplicative hazard has a clean reading: the coefficients $\gamma_1,\gamma_2,\gamma_3$ act as behavioral elasticities (when $\gamma_2\approx\ln 2$, doubling loyalty roughly halves the hazard), and they can differ across cohorts---first-year students often have high $\gamma_1$ but low $\gamma_2$ (enthusiasm over habit), graduating students the reverse. Linearizing for small hazards, dropout probability $\approx h_i$ and reactivation probability $\approx\xi_1 R_i+\xi_2\bar a_{N(i)}$, clarifying sensitivities. Because dropout thins edges and reactivation restores them, the network ``breathes'': its effective density contracts and expands with membership. In steady state the active fraction is approximately $\mathrm{react}/(\mathrm{react}+\mathrm{drop})$, so participation is cyclical rather than monotone---an open demographic ecosystem in which the collective rhythm persists through the interplay of decay and renewal. As long as $\xi_1,\xi_2>0$ and $h_{0,i}>0$, the active/inactive chain is ergodic, yielding measurable quantities such as expected tenure, turnover rate, and mean reactivation delay; the parameters are estimable from longitudinal attendance via survival analysis.

\subsection{Numerical stability}\label{sec:stability}
The soft-saturation transform $\tilde X_i=X_i/(1+\rho_i X_i)$ in \eqref{eq:softsat} performs three roles at once. Mathematically, it keeps the coupled iteration contractive: as $X_i\to\infty$ it approaches $1/\rho_i$, and increasing $\rho_i$ lowers the slope of the peer-response function and hence the spectral radius of the linearized interaction map, preventing runaway growth when many motivated nodes reinforce one another. Computationally, unlike hard truncation it remains differentiable, enabling gradient-based sensitivity analysis and reverse planning. Behaviorally, it encodes \emph{limited receptivity}: students have finite emotional bandwidth, so additional peer signal yields diminishing returns and $\rho_i$ is interpretable as ``social resistance.'' A sufficient global-boundedness condition is
\begin{equation*}
\rho_i > w^{P}_i\textstyle\sum_j \omega_{ij}(t)\,b_j\quad\text{for all }i,
\end{equation*}
which guarantees the maximum possible influence input never exceeds the saturation threshold; it is cheap to check and can be embedded as a safety constraint. For continuous-time analogues, explicit Euler with step $\Delta t<1/\max_i\alpha_{0,i}$ ensures stability, and all states are projected to $[0,1]$ after each update.

\subsection{Modulating states: emotion, saturation, stability}\label{sec:modulating}
Three slower states modulate the core dynamics and are given explicit updates. Emotion mean-reverts to a temperament baseline and is nudged by the valence of reinforcement:
\begin{equation}\tag{E}\label{eq:emotion}
E_i(t{+}1)=\clip\!\Big[(1-\lambda_E)E_i(t)+\lambda_E\bar E_i+\zeta_E[B_i(t)]_+-\zeta'_E[B_i(t)]_-+\kappa_E\,\xi^E_i(t)\Big].
\end{equation}
Social saturation slowly tracks the normalized incoming-weight load, gated by activity:
\begin{equation}\tag{S}\label{eq:saturation}
\ell_i(t)=\sum_j\omega_{ij}(t);\qquad S_i(t{+}1)=\mathrm{act}_i\cdot\Big[(1-\lambda_S)S_i(t)+\lambda_S\,\tfrac{\ell_i(t)}{\ell_i(t)+s_0}\Big].
\end{equation}
Stability falls with recent motivation variance and enters \eqref{eq:mot}--\eqref{eq:fatigue} through the gain $G(\Phi)$:
\begin{equation}\tag{$\Phi$}\label{eq:stability}
\begin{aligned}
\bar m_i(t)&=(1-\lambda_v)\bar m_i(t{-}1)+\lambda_v m_i(t),\\
\mathrm{var}_i(t)&=(1-\lambda_v)\mathrm{var}_i(t{-}1)+\lambda_v\big(m_i(t)-\bar m_i(t)\big)^2,\\
\Phi_i(t)&=\frac{1}{1+\kappa_v\,\mathrm{var}_i(t)}.
\end{aligned}
\end{equation}
Their rates $\lambda_E,\lambda_S,\lambda_v$ are slow (much smaller than $\lambda_L$) and the coupling gains $\kappa_\bullet$ are small. Setting the modulating gains $\beta_E,\kappa_{EF},\kappa_{SL},\kappa_\Phi$ to zero and the saturation/mood multipliers to $1$ recovers the core, un-modulated model exactly, so the three states are clean optional refinements: $E$ lets affect amplify or dampen reinforcement, $S$ encodes diminishing peer returns under crowding, and $\Phi$ stabilizes erratic trajectories.

\subsection{Modeling assumptions and functional-form rationale}\label{sec:assumptions}
Each functional form encodes a specific, falsifiable behavioral hypothesis; the framework is modular, and the forms may be substituted without changing its structure provided boundedness and the qualitative monotonicities are preserved. I record the rationale explicitly so that alternatives can be tested rather than assumed away.

The \emph{multiplicative} decay \eqref{eq:decay} is chosen over an additive combination of loyalty and fatigue effects because the two should \emph{interact} rather than merely sum: a highly loyal but exhausted member should not decay at a rate obtained by adding a large protective term to a large depleting one, but at a rate where exhaustion erodes the protection loyalty would otherwise provide. The product form also keeps $\alpha_i\in(0,1)$ automatically and makes loyalty's protective effect proportionally larger exactly when the baseline decay $\alpha_{0,i}$ is high, which is where protection matters most. The \emph{concave-with-sign-flip} internal reinforcement \eqref{eq:internal} bundles three hypotheses: diminishing returns to effort, through the concave $\beta_i(1-e^{-\delta_i y_i})$; a qualitative rather than graded distinction between minimal participation and none, through the discrete penalty $-\chi_i\mathbf 1\{y_i=0\}$; and the discouragement of an under-attended session, through the quorum sign flip that turns reinforcement flatly negative below $N_{\min}$ regardless of individual preparedness. This last feature is the one most responsible for domain-dependent collapse. The \emph{proportional-hazards} dropout form \eqref{eq:hazard} is multiplicative in $m_i,L_i,F_i$ so that each coefficient reads as a behavioral elasticity, guarantees $h_i>0$, and nests the constant-hazard (memoryless tenure) model when $\gamma_1=\gamma_2=\gamma_3=0$. The \emph{rational} soft-saturation \eqref{eq:softsat} is preferred over hard truncation because it is everywhere differentiable---a prerequisite for the gradient-based planning of Sections~\ref{sec:control}--\ref{sec:feasibility}---and admits the clean contraction bound of Section~\ref{sec:stability}; the $\tanh$ transmission $g$ bounds each neighbor's contribution so that no single enthusiast dominates the peer sum. Finally, the \emph{logistic} action map \eqref{eq:actions} represents a participation decision as a noisy threshold whose slope $\kappa$ interpolates between a deterministic threshold ($\kappa\to\infty$) and a near-linear probability ($\kappa\to0$), with separate slopes for attending and posting capturing that the two decisions respond differently to the same underlying motivation.

\subsection{Fixed points, bistability, and the collapse--stabilization dichotomy}\label{sec:fixedpoints}
A mean-field reduction makes the qualitative behavior transparent and shows that the contrast between a stabilizing and a collapsing club is not an artifact of simulation but a property of the fixed-point structure of \eqref{eq:mot}--\eqref{eq:internal}. Suppress heterogeneity and noise, set the stability gain $G(\Phi)=1$, take emotion at baseline so that $\hat m_i=m_i$ and $\nu_E=1$, and seek a symmetric steady state in which every node carries the same triple $(m,L,F)$ on a regular network. Write the common attendance probability as
\[
q(m)=\sigma\!\big(\kappa^a(m-\theta^a)\big),
\]
and assume operation below capacity, so that effective attendance $\tilde a\approx q(m)$ and the neighbor average $\bar a_{N(i)}\approx q(m)$. Solving the loyalty and fatigue updates \eqref{eq:loyalty}--\eqref{eq:fatigue} at steady state expresses both stocks as increasing functions of participation,
\[
L^\star(m)=(\varphi_1+\varphi_2 s^{\mathrm{role}})\,q(m)+\varphi_3 r^+,
\qquad
F^\star(m)=(\psi_1\bar y+\psi_2-\psi_3)\,q(m)+\psi_3,
\]
with $y\approx\bar y\,q(m)$. Substituting into the decay law \eqref{eq:decay} (with $S$ small) gives the effective decay as a function of participation,
\[
\alpha^\star(m)=\alpha_{0}\big(1-\eta_L L^\star(m)\big)\big(1+\eta_F F^\star(m)\big),
\]
which \emph{falls} as $m$ rises whenever loyalty's protective slope outweighs fatigue's---the central stabilizing feedback of the model. Collecting the reinforcing drive from peers, events, and internal progress,
\[
D(m)=w^{P}\,\tilde X\big(q(m)\big)+c\,R+B\big(q(m)\big),
\]
where $B(\cdot)$ from \eqref{eq:internal} is positive and saturating above quorum and flatly negative below it, the symmetric steady state is any solution of the scalar self-consistency equation
\begin{equation}\tag{$\star$}\label{eq:selfconsist}
\alpha^\star(m)\,m=D(m)
\qquad\Longleftrightarrow\qquad
m=\frac{D(m)}{\alpha^\star(m)}=:T(m).
\end{equation}
Because $q(m)$ is sigmoidal, $D(m)$ is small for low $m$, rises steeply through the threshold region, and saturates, while $\alpha^\star(m)$ is high at low $m$ (no loyalty, elevated fatigue) and lower at high $m$. The composite map $T(m)$ is therefore itself S-shaped and can cross the identity line at up to three points: a low fixed point $m_0$ (a \emph{collapsed} club---few attend, no loyalty accumulates, decay stays high), an unstable middle point $m_u$ (the tipping threshold), and a high-participation fixed point $m_1$ (an \emph{active} club---many attend, loyalty has accrued, decay is suppressed). The system is bistable.

This single picture organizes the rest of the paper. The \emph{stabilization} regime corresponds to parameters---low $\alpha_0$, large $\beta$, small $\delta$, low $N_{\min}$, high $w^P$ and $V^{ex}$---under which $D(m)$ is large and $\alpha^\star(m)$ small enough that the high-participation fixed point $m_1$ exists with a wide basin of attraction; the \emph{collapse} regime is the opposite, with the quorum holding $B$ negative across the relevant range and $\alpha^\star$ high, so that the high branch is destroyed and only $m_0$ survives. The transition between the two as parameters vary is a fold (saddle-node) bifurcation in which the high-participation fixed point and the tipping threshold collide and annihilate. The feasibility frontier of Section~\ref{sec:feasibility} is exactly this fold---the parameter boundary at which a self-sustaining state ceases to exist---and a participation target above it is structurally infeasible because no bounded input can stabilize a state the dynamics do not possess. Interventions read cleanly in these terms: a one-shot event raises the constant $R$, lifting $D(m)$ and pushing the state above $m_u$, but if the regime admits only the collapsed fixed point the high state is a transient and the trajectory returns to $m_0$ once $R$ decays. A successful intervention instead changes the fixed-point structure itself, by holding participation high long enough for $L^\star(m)$ to accumulate, thereby lowering $\alpha^\star(m)$ and bringing the high-participation fixed point into existence (or widening its basin), so that withdrawing the input leaves the system resting on $m_1$ rather than $m_0$. This is the precise sense in which control must ``move the slow variable,'' and it is why temporary scaffolding can produce a permanent change in outcome. Heterogeneity and real network structure enrich the picture---thresholds vary across students, and clustering localizes the tipping behavior---but the bistable skeleton survives.

\subsection{Limiting cases, observation, and interpretation}\label{sec:nesting}
The observable variables are $\{a_i(t),p_i(t)\}$; the states $(m_i,L_i,F_i)$ are latent, and $(E_i,S_i,\Phi_i)$ evolve deterministically from observable histories, entering estimation only through their effects on $m_i,\alpha_i,\omega_{ij}$. The framework nests the classical models: holding $L,F,B$ constant, freezing $\omega$, and taking $\kappa\to\infty$ reduces \eqref{eq:mot} to a Linear Threshold cascade, while a single one-shot stochastic pulse reproduces Independent Cascade behavior. The additional states govern departures from memoryless diffusion toward habit formation, saturation, recovery, and turnover. The system differs from product diffusion in three respects: continuous motivation rather than binary adoption; bidirectional feedback between personal experience and network exposure; and a heterogeneous, evolving influence structure. These properties make it suitable for both forward simulation and the reverse planning developed next.

\paragraph{Per-week update order.}
The dynamics are \emph{synchronous}: within each week, all time-$t$ observables and inputs (the actions and effective attendance \eqref{eq:actions},\eqref{eq:capacity}; the time-$t$ weights and aggregates \eqref{eq:omega}--\eqref{eq:cred}; the soft-saturated peer input \eqref{eq:softsat}; the reinforcements \eqref{eq:external},\eqref{eq:internal}; the effective fatigue, decay, and gain \eqref{eq:decay}; and the stability factor \eqref{eq:stability}) are evaluated from the current state \emph{before} any state is advanced to $t{+}1$ via \eqref{eq:mot},\eqref{eq:loyalty},\eqref{eq:fatigue},\eqref{eq:emotion},\eqref{eq:saturation},\eqref{eq:vupdate}, followed by attrition/reactivation \eqref{eq:hazard},\eqref{eq:react}. In implementation this is a read-from-current, write-to-next double buffer; the full ordering and a reference implementation accompany the paper as supplementary material.

\subsection{Identification and estimation}\label{sec:estimation}
The observable variables are the actions $\{a_i(t),p_i(t)\}$; the core states $(m_i,L_i,F_i)$ are latent, and the modulating states $(E_i,S_i,\Phi_i)$ evolve deterministically from observable histories. Estimation therefore proceeds at two levels. Given member-level attendance and posting logs, the probabilistic action map \eqref{eq:actions} defines a state-space model whose parameters can be fit by a particle-filter expectation--maximization scheme: the filter infers the latent trajectory $(\hat m_i,\hat L_i,\hat F_i)_t$ from observed actions, and the maximization step updates the structural parameters, with spells of inactivity treated as right-censored tenure rather than as zero attendance. When only group-level aggregates are available---attendance counts, posting volume, membership---likelihood-free methods such as Approximate Bayesian Computation or simulated minimum distance match summary trajectories (expected attendance, turnover rate, mean tenure) between simulation and data.

Identifiability deserves care, because several parameters trade off and require informative variation to separate. The baseline decay $\alpha_{0,i}$ and the loyalty sensitivity $\eta_L$ both govern how fast motivation erodes and are distinguishable only when loyalty itself varies across the observation window. The external channel $R_i$ \eqref{eq:external} and the internal channel $B_i$ \eqref{eq:internal} both raise motivation around events and are separated only by exploiting event timing relative to participation: $R_i$ acts immediately and decays exponentially, whereas $B_i$ depends on realized attendance and the quorum. The hazard elasticities $\gamma_1,\gamma_2,\gamma_3$ \eqref{eq:hazard} are recoverable from the covariation of dropout with the inferred states, which argues for collecting enough longitudinal dropout events to fit a proportional-hazards model. These are not obstacles but design requirements: event variation, a sufficient observation horizon, and member-level rather than only aggregate logs are precisely what make the model identifiable.

% =====================================================================
\section{Reverse planning and cost-minimizing control}\label{sec:control}

\subsection{From forward dynamics to inverse design}
The forward system generates trajectories from initial conditions, domain parameters, and intervention schedules. But organizers rarely ask what will happen under a fixed configuration; they ask how a desired outcome can be achieved under resource, capacity, and organizational constraints. That problem is inverse: infer interventions that steer the system to a target while minimizing cost. I reformulate the model as a controllable dynamical system in which all behavioral dynamics remain governed by \eqref{eq:actions}--\eqref{eq:react}; no new behavioral assumptions are added, and the equations are treated as constraints. Time is weekly with horizon $T=18$ (a semester), and $V$ is the whole school, interpreted as the set of recruitable students, on a network with moderate clustering and a non-trivial fraction of bridging ties.

\subsection{Control variables and feasible intervention space}
Controls $u$ enter only through existing channels: (i) the event schedule $\{(\tau_k,A_k)\}$ via the pulse \eqref{eq:external}, bounded $0\le R_i\le\bar R$ with at most $M$ events; (ii) visibility amplification $V^{ex}_i\leftarrow\gamma_v V^{ex}_i$ via \eqref{eq:vmix}--\eqref{eq:vupdate}, with $1\le\gamma_v\le\bar\gamma_v$; (iii) the capacity schedule $K(t)$ via \eqref{eq:capacity}; and (iv) early scaffolding $\gamma_B(t)$ on the internal reinforcement \eqref{eq:internal} over a short window. Each lever has a natural cost (staff time, space, attention budget), and each acts on a different part of the mechanism: events and scaffolding move the fast variable, while visibility and capacity reshape the channels through which the slow variables accumulate.

\subsection{Objective and expected attendance}
The planner minimizes intervention cost subject to the dynamics and a terminal-window participation target:
\begin{equation*}
\min_{u}\ C(u)\quad\text{s.t.\ Eqs.\,\eqref{eq:actions}--\eqref{eq:react} and}\quad
\bar A_{14:18}(u)=\tfrac15\!\!\sum_{t=14}^{18}\!A(t)\ \ge\ A^\star,\qquad A(t)=\sum_i\tilde a_i(t).
\end{equation*}
Because actions are probabilistic, expected attendance is computed analytically by summing individual probabilities rather than by Monte Carlo alone:
\begin{equation*}
\mathbb{E}[A(t)]=\sum_i P^a_i(t)\,\mathrm{act}_i.
\end{equation*}
This expectation-based aggregation is what makes both the objective and the feasibility screen of Section~\ref{sec:feasibility} cheap to evaluate.

The cost is additive across levers and weeks,
\begin{equation*}
C(u)=\sum_t\Big[\kappa_E\,n_E(t)+\kappa_v\big(\gamma_v(t)-1\big)+\kappa_K\,K(t)+\kappa_B\big(\gamma_B(t)-1\big)\Big],
\end{equation*}
where $n_E(t)$ counts events scheduled at week $t$ and each coefficient converts a lever into its real resource cost---staff hours to run an event, attention budget to amplify visibility, space and supervision to expand capacity, and mentoring effort to scaffold early reinforcement. Because the levers act on different parts of the mechanism---events and scaffolding move the fast variable $m$ directly, while visibility and capacity reshape the channels through which loyalty and the influence weights accumulate---the optimal schedule is rarely a single lever, and the planner trades immediate, decaying boosts against slower but persistent structural change.

\subsection{Domain regimes}\label{sec:regimes}
Parameters are drawn per student from regime distributions (mean\,$\pm$\,sd); one period is one week. Table~\ref{tab:regimes} contrasts a \emph{collapse} (composition-like) regime against a \emph{stabilize} (fitness-like) regime. The collapse regime represents high skill, delayed feedback, low expressive visibility, and a strong quorum effect; the stabilization regime represents low entry barriers, immediate feedback, high visibility, and productivity even at small group sizes.

\begin{table}[t]
\centering\small
\setlength{\tabcolsep}{8pt}
\renewcommand{\arraystretch}{1.12}
\begin{tabular}{@{}lcc@{}}
\toprule
Parameter & Collapse (composition) & Stabilize (fitness)\\
\midrule
$m_i(0)$ initial motivation & $0.28\pm0.10$ & $0.33\pm0.10$\\
$L_i(0)$ initial loyalty & $0.10\pm0.05$ & $0.12\pm0.06$\\
$F_i(0)$ initial fatigue & $0.35\pm0.10$ & $0.28\pm0.10$\\
$\alpha_{0,i}$ baseline decay & $0.45\pm0.10$ & $0.28\pm0.08$\\
$\eta_L$ decay\,$\leftarrow$\,loyalty & $0.70$ & $0.75$\\
$\eta_F$ decay\,$\leftarrow$\,fatigue & $0.55$ & $0.35\pm0.12$\\
$w^{P}_i$ peer susceptibility & $0.25\pm0.10$ & $0.45\pm0.20$\\
$b_i$ spreading propensity & $0.22\pm0.12$ & $0.48\pm0.10$\\
$\theta^a_i$ attendance threshold & $0.62\pm0.10$ & $0.58\pm0.12$\\
$\theta^p_i$ posting threshold & $0.72\pm0.10$ & $0.72\pm0.15$\\
$\kappa^a_i,\kappa^p_i$ slopes & $9,\ 11$ & $8,\ 9$\\
$D_i$ domain preparedness & $0.30\pm0.15$ & $0.70\pm0.15$\\
$V^{ex}_i$ expressive visibility & $0.25\pm0.12$ & $0.60\pm0.15$\\
$c_i$ availability & $0.55\pm0.15$ & $0.55\pm0.15$\\
$E_i(0)$ emotion & $0.45\pm0.15$ & $0.22\pm0.07$\\
$\beta_i$ internal-reinf.\ cap & $0.12\pm0.06$ & $1.4\pm0.4$\\
$\delta_i$ reinf.\ sensitivity & $0.9\pm0.3$ & $0.06\pm0.03$\\
$\chi_i$ skip penalty & $0.10\pm0.05$ & regime default\\
$N_{\min}$ quorum & $10$ & $4$\\
\bottomrule
\end{tabular}
\caption{Two domain regimes (per-student draws). The sign-flip quorum $N_{\min}$ and the reinforcement pair $(\beta_i,\delta_i)$ are the decisive differences.}
\label{tab:regimes}
\end{table}

\subsection{Three reference scenarios}\label{sec:scenarios}
Within each case the regime is fixed; differences across scenarios come solely from the control vector (Table~\ref{tab:scenarios}). Attendance values are expected totals. Scenario A applies a single large event in a collapse regime and decays below quorum. Scenario B applies modest, repeated events with amplified visibility and ample capacity in a stabilizing regime, settling at a fixed point above target. Scenario C is a reverse-planned rescue of a collapse-regime club: a small event, strong visibility amplification, and a short scaffolding window hold participation above $N_{\min}$ long enough for loyalty to accumulate and cut the decay rate.

\begin{table}[t]
\centering\small
\setlength{\tabcolsep}{6pt}
\renewcommand{\arraystretch}{1.15}
\begin{tabular}{@{}lp{3.4cm}p{3.4cm}p{3.6cm}@{}}
\toprule
 & A. Naive (collapse) & B. Favorable (stable) & C. Reverse-planned rescue\\
\midrule
regime & collapse & stabilize & collapse\\
events $(\tau,A)$ & one, $(2,0.60)$ & $(2,6,10)$ at $A=0.35$ & one, $(2,0.30)$\\
$\gamma_v$ & $1$ & $1.4$ & $3.5$\\
$K(t)$ & $14$ & $18$ & $14$\\
$\gamma_B(t)$ & $1$ & $1$ & $1.8$ on $t\in\{2..6\}$\\
$\mathbb{E}[A]$ trajectory & wk 3$\to$18: 9.1, 4.0, 2.6, 1.7, 0.9 & wk 3$\to$18: 13.8, 14.0, 14.0, 14.0, 14.0 & wk 10$\to$18: 8.3, 7.9, 7.4, 6.8, 5.7\\
outcome & decays below $N_{\min}$, collapses & stable point above target & above $N_{\min}$ early; loyalty accrues; survives\\
\bottomrule
\end{tabular}
\caption{Three reference scenarios. Differences across columns are entirely in the control vector, not the regime.}
\label{tab:scenarios}
\end{table}

\subsection{Mechanism-level walkthrough}\label{sec:walkthrough}
The three scenarios are most informative read against the bistability of Section~\ref{sec:fixedpoints}, since each is a different relationship between the intervention and the fixed-point structure.

\emph{Scenario A (naive collapse).} The single large event spikes $R_i$ and pushes attendance above the tipping threshold for the first two weeks, so early counts are high. But the regime is collapse: with $N_{\min}=10$ and small $\beta_i$, realized attendance falls below quorum as soon as the pulse begins to decay, the internal term $B_i$ \eqref{eq:internal} turns negative, and because no loyalty has had time to accumulate the effective decay $\alpha^\star$ remains near its baseline of $0.45$. Motivation reverts to the only fixed point the regime admits, and attendance decays monotonically from $9.1$ toward $0.9$. The event moved the fast variable and nothing moved the slow one.

\emph{Scenario B (favorable, self-sustaining).} Modest repeated events keep $R_i$ topped up while the stabilization regime already admits a wide-basin high-participation fixed point. Capacity $K=18$ exceeds demand, so effective attendance is never clipped and loyalty and credibility accumulate freely; rising loyalty lowers $\alpha^\star$, and the club settles onto the high-participation fixed point with attendance plateauing at $14.0$. The defining feature is that the plateau persists as events space out---the system rests at $m_1$ rather than depending on continual stimulation.

\emph{Scenario C (reverse-planned rescue).} The regime is the same collapse regime as Scenario A, which initially admits no high-participation fixed point. The rescue combines a small event, strong visibility amplification $\gamma_v=3.5$ (which thickens edges through \eqref{eq:vupdate} and raises the peer term), and a short scaffolding window $\gamma_B=1.8$ over weeks $2$--$6$ (which makes the otherwise-negative early $B_i$ positive while participation is still small). Together these hold participation above $N_{\min}$ through the window; during that window loyalty climbs and $\alpha^\star$ falls, which brings a high-participation fixed point into existence. When the scaffolding ends at week~$6$ the now-reduced decay sustains participation: attendance decays only gently, from $8.3$ toward $5.7$, and stays above quorum. A temporary input has produced a permanent change in the fixed-point structure---the operational signature of controllability.

\subsection{Interpretation}
Participation diffusion is structurally controllable. Collapse arises when interventions act only on fast variables such as motivation: a one-off pulse raises attendance briefly, but without moving the slow stabilizers the trajectory returns to the decay-dominated basin. Success emerges when interventions shift the slow variables---loyalty, and hence the effective decay---by holding participation above $N_{\min}$ long enough for reinforcement to compound. Reverse planning identifies when rescue is feasible and when failure reflects a domain mismatch rather than insufficient effort.

% =====================================================================
\section{Feasibility, closed-loop control, and constrained design}\label{sec:feasibility}

\subsection{Feasibility frontier}
A controllable system is still a constrained one. With event pulses bounded in amplitude and count, visibility bounded by $\bar\gamma_v$, and realized participation capped by $K(t)$, the achievable set of terminal-window participation has a supremum, and a target is attainable only if it lies below it:
\begin{equation*}
A^\star\ \text{attainable}\iff A^\star\le\sup_{u\in\mathcal U}\bar A_{14:18}(u).
\end{equation*}
This separates two failure modes that look identical in raw outcomes: \emph{optimization failure}, and \emph{structural infeasibility}. The latter arises when bounded controls cannot push the system out of the decay-dominated basin---characterized by persistently high $\alpha_i$ \eqref{eq:decay}, weak or negative $B_i$ \eqref{eq:internal}, and elevated hazard \eqref{eq:hazard}---so that early gains from $R_i$ dissipate faster than loyalty can accumulate. A fast diagnostic maximizes $\sum_i P^a_i(t)$ week-by-week under the bounds and checks whether the resulting trajectory stays above the minimum-viability region long enough for loyalty to reduce decay. This screen does not replace simulation; it answers immediately whether a target is impossible under current ceilings even with a perfectly designed intervention. Capacity bounds make the point sharp: once demand exceeds $K(t)$, raising $\bar R$ or $\bar\gamma_v$ yields diminishing returns because clipped participation slows the accumulation of loyalty and credibility. Reverse planning is therefore also a tool for proving negative results: if $A^\star$ lies above the frontier, no schedule delivers it without changing the domain regime itself.

\subsection{Closed-loop control via receding-horizon re-optimization}
Deployment is sequential: the organizer observes attendance and posting weekly while latent states and the influence network $W_t$ evolve in the background. The probabilistic action map \eqref{eq:actions} permits filtering---at the end of week $t$, observations $y_t$ yield estimates $(\hat x_t,\hat W_t)$ with $\hat x_t=\{\hat m_i,\hat L_i,\hat F_i\}$ from, e.g., a particle filter. Closed-loop planning then solves a receding-horizon problem, applies only the first control, and repeats:
\begin{equation*}
u^\star_{t:t+H-1}=\argmin_{u_{t:t+H-1}\in\mathcal U}\ \mathbb{E}\Big[\textstyle\sum_{\tau=t}^{t+H-1}c(u_\tau)\Big]
\quad\text{s.t.\ Eqs.\,\eqref{eq:actions}--\eqref{eq:react},}\ (\hat x_t,\hat W_t),\ \text{terminal feasibility.}
\end{equation*}
This has three advantages absent in open-loop form. It adapts to \emph{parameter drift}---exam weeks shift fatigue, the hazard \eqref{eq:hazard}, and reactivation \eqref{eq:react}---because each week's estimate revises the trajectory. It \emph{prevents over-intervention}: once loyalty accumulates and $\alpha_i$ \eqref{eq:decay} declines, the scheme detects the onset of self-stabilization and withdraws effort without risking collapse, operationalizing the claim that the system is controllable rather than merely boostable. And it provides a direct bridge to an implementable interface that displays predicted attendance, inferred loyalty and fatigue, and a recommended control that is the first action of an optimized plan rather than a heuristic. Churn is handled by stating the participation constraint in terms of expected realized attendance after dropout and capacity.

\subsection{Constrained optimization under load, equity, and robustness}
A controllable system invites concentrating interventions on the most influential nodes---an ethical and organizational risk, since the model itself converts overuse into harm through fatigue and the dropout hazard \eqref{eq:hazard}. The following constraints add no behavioral structure; they modify only the feasible set.

\emph{Individual load.} Let $\ell_i(t)$ be the intervention-induced load on student $i$, e.g.\ $\ell_i(t)=\omega_R R_i(t)+\omega_v\,\mathrm{outgoing\_visibility}_i(t)+\omega_p\,\mathrm{posting\_targeting}_i(t)$, each term computable from the model. The fairness constraint is $\sum_t \ell_i(t)\le\Lambda_i$ for all $i$ (or for a protected subset such as already-fatigued or highly central students). Enforcing it directly is better than relying on the hazard to punish unrealistic plans after the fact.

\emph{Equity across subgroups.} Because \eqref{eq:vupdate} strengthens ties through co-activity, visibility amplification can make strong clusters stronger and deepen inequality. With subgroup terminal-window averages $\bar A_g=\tfrac15\sum_{t=14}^{18}\sum_{i\in g}\tilde a_i(t)$, an equity constraint $\bar A_g\ge\rho_g A^\star$ prevents the optimizer from hitting the global target by sacrificing peripheral groups, without forcing identical outcomes.

\emph{Robustness under parameter uncertainty.} Baseline parameters are plausible rather than measured for every individual. With uncertain parameters $\theta\in\Theta$ (e.g.\ $\beta_i,\alpha_{0,i}$, and the hazard elasticities), robust planning solves a worst-case problem:
\begin{equation*}
\min_{u\in\mathcal U}\ \sup_{\theta\in\Theta}\ C(u;\theta)
\qquad\text{s.t.}\qquad \inf_{\theta\in\Theta}\bar A_{14:18}(u;\theta)\ge A^\star.
\end{equation*}
This matters most in the rescue regime, which operates near the frontier, where small misestimation of decay or fatigue sensitivity can flip the sign of reinforcement persistence and send the system back toward collapse. Together these constraints change what ``optimal'' means: the cheapest unconstrained plan may lean on a few credible hubs, whereas the constrained plan diversifies influence pathways, increases the role of network shaping through \eqref{eq:vupdate}, and turns events from pure boosts into coordination tools that distribute activation across the network while respecting load limits.

\subsection{From plan to practice: an organizer-facing workflow}\label{sec:workflow}
The closed-loop formulation maps onto a deployable workflow that does not require the organizer to understand the model. Each week the organizer records who attended and who posted; these observations drive the filter of Section~\ref{sec:feasibility}, which updates the inferred motivation, loyalty, and fatigue of each member and the current influence network. The receding-horizon problem is then re-solved, and the organizer is shown three things: a short-horizon forecast of expected attendance, a summary of which members are accumulating loyalty versus drifting toward dropout, and a single recommended action for the coming week---the first control of the optimized plan---together with its predicted effect. The recommendation is an optimized first step rather than a heuristic, and the same machinery that recommends intensification also recommends \emph{withdrawal}: once loyalty has accumulated and $\alpha^\star$ has fallen, the controller detects the onset of self-stabilization and reduces effort, which both conserves the organizer's limited resources and respects the load and equity constraints above. This withdrawal behavior is what distinguishes a controllable system from one that must be continually propped up, and it is the practical payoff of the bistability analysis: the aim is not to maximize stimulation but to deliver the system into the basin of the high-participation fixed point and then let it rest there.



\subsection{Interactive companion and ensemble readouts}\label{sec:companion}
The same equations can be exposed through an organizer-facing simulator without modifying the model. The companion interface therefore separates two computational views. A \emph{sample} run draws one randomized influence network and one set of heterogeneous initial conditions from the specified regime, applies a chosen control path $u$, and displays the full weekly trajectory with a time scrubber. This view is useful for mechanism-level interpretation: one can see which core students cross the attendance threshold, which students accumulate loyalty, and where dropout occurs.

A \emph{simulation} run instead treats the graph generator and heterogeneous parameter draws as an ensemble distribution. For $M$ independent realizations indexed by $r$, with the same control path $u$ and horizon $T$, the displayed attendance curve is
\begin{equation}
  \widehat A_M(t;u)=\frac{1}{M}\sum_{r=1}^{M} A^{(r)}(t;u),
  \qquad t=0,\dots,T,
  \label{eq:ensembleA}
\end{equation}
and analogous averages are shown for mean motivation, loyalty, fatigue, active membership, and reach. Percentile bands, such as the interquartile interval of $\{A^{(r)}(t;u)\}_{r=1}^M$, are visualization summaries of network sensitivity rather than new state variables. Node-level averages in the ensemble display should be read as averages over indexed simulated agents, or over comparable positions in the generated core network, not as predictions about named students.

The close-ties-per-student and bridge-ties-per-student controls parameterize the distribution from which the directed network is sampled. They change the graph ensemble, not the behavioral update laws \eqref{eq:actions}--\eqref{eq:react}. Likewise, the public-body dots in the visualization are a presentation layer representing students outside the simulated core network: they help the viewer distinguish a core influence network from the wider school population, but they do not enter the equations unless explicitly promoted into the model as additional nodes in a future extension. In reverse mode, the optimizer first returns a feasible proposed control path $u^\star$; the sample-plan and simulation-plan controls simply evaluate that same $u^\star$ under the two readout modes above. Thus the interface adds randomized robustness checks and explanatory visualization while preserving the original model as the governing dynamical system.

% =====================================================================
\section{Discussion and limitations}\label{sec:discussion}
The model trades the analytic cleanliness of binary, static-graph diffusion for behavioral realism: continuous and renewable participation, endogenous decay, co-evolving influence, and explicit turnover. This realism is what lets a single framework reproduce both stabilization and collapse from the same equations under different regimes, and what makes intervention design a constrained-optimization problem with a feasibility frontier rather than a search for the largest seed set. The interactive companion should be read in the same spirit: its sample and ensemble modes are numerical illustrations of the equations, not calibrated forecasts of a particular school. Several limitations bound the claims. The parameters are plausible and illustrative rather than estimated for a specific population; the reference trajectories in Section~\ref{sec:scenarios} demonstrate qualitative regimes, not calibrated forecasts, and reproducing them quantitatively requires fitting the model to longitudinal attendance and posting logs. The behavioral laws---the multiplicative decay \eqref{eq:decay}, the concave-with-sign-flip reinforcement \eqref{eq:internal}, and the exponential hazard \eqref{eq:hazard}---are modeling choices motivated by the literature and should be tested against alternatives. Finally, the control results assume an estimator good enough to recover latent states from sparse observations; the quality of closed-loop control degrades with observation noise, and the responsible-design constraints of Section~\ref{sec:feasibility} are necessary precisely because an unconstrained optimum can be ethically inappropriate in a youth setting. These are directions for empirical calibration and validation rather than obstacles to the framework's structure.

The model is falsifiable in specific ways. It predicts bistability and hysteresis: clubs with similar current participation but different histories should diverge, and a club pushed above its tipping threshold and then released should either remain active or fall back depending on whether loyalty crossed the level that lowers decay. It predicts that the same event produces persistent gains in low-barrier domains and merely transient gains in quorum-dependent ones. And it predicts that reinforcement, not exposure alone, governs survival, so that interventions raising visibility without raising realized, above-quorum participation should fail. These are testable against the calibration data described below. A final caution is ethical rather than empirical: because the model converts overuse of influential members into fatigue and dropout, and because visibility amplification can deepen inequality between subgroups, the responsible-design constraints of Section~\ref{sec:feasibility} are not optional refinements but conditions for using the framework appropriately in a youth setting, and any deployment should be transparent about what is being optimized and should keep a human organizer in the loop.

% =====================================================================
\section{Future work and empirical calibration}\label{sec:future}
The framework is built to be tested, and the natural next step is calibration against longitudinal data from real clubs. The data requirements follow from the identification analysis of Section~\ref{sec:estimation}: member-level attendance and posting records over a full term, timestamps and magnitudes of organized events, and---where feasible---light periodic measures of affect or perceived progress to anchor the emotion and internal-reinforcement channels. With such data, the estimation pipeline infers latent trajectories and structural parameters, and the fitted model can be validated in three complementary ways: by held-out trajectory prediction (fit on early weeks, predict later participation); by posterior predictive checks on the full distribution of weekly attendance rather than only its mean; and by ablation---removing each modulating state in turn and measuring the loss in fit---to establish which mechanisms the data actually require.

Several extensions are natural and each preserves the core dynamics while adding a channel. Clubs do not exist in isolation: students allocate a shared budget of time and attention across competing activities, so coupling several instances of the model through common capacity and availability would turn the single-club analysis into a study of competition for participation, in which one club's growth can suppress another's. The event model can be enriched with heterogeneous event types---recruitment drives, competitions, social events---that load the reinforcement and visibility channels differently. Affect, treated here as an individual state, could be made contagious across edges so that collective mood spreads alongside participation. Seasonality such as exam periods and holidays enters cleanly as time-varying fatigue and availability. Finally, leadership is currently exogenous; endogenizing the emergence of high-credibility members would connect the model to the organizational question of how clubs develop the internal structure that sustains them. The intent throughout is a framework that is modular rather than final---one in which each additional mechanism is an explicit, testable channel rather than a hidden assumption.

% =====================================================================
\section*{References}
\begin{flushleft}\small
Aral, S., \& Walker, D. (2012). Identifying influential and susceptible members of social networks. \textit{Science, 337}, 337--341.\\[2pt]
Bond, R. M., et al. (2012). A 61-million-person experiment in social influence and political mobilization. \textit{Nature, 489}, 295--298.\\[2pt]
Centola, D. (2010). The spread of behavior in an online social network experiment. \textit{Science, 329}, 1194--1197.\\[2pt]
Granovetter, M. (1978). Threshold models of collective behavior. \textit{American Journal of Sociology, 83}, 1420--1443.\\[2pt]
Gross, T., \& Blasius, B. (2008). Adaptive coevolutionary networks: A review. \textit{J. R. Soc. Interface, 5}, 259--271.\\[2pt]
Holme, P., \& Newman, M. E. J. (2006). Nonequilibrium phase transition in the coevolution of networks and opinions. \textit{Physical Review E, 74}, 056108.\\[2pt]
Holme, P., \& Saram\"aki, J. (2012). Temporal networks. \textit{Physics Reports, 519}, 97--125.\\[2pt]
Kairam, S. R., Wang, D. J., \& Leskovec, J. (2012). The life and death of online groups. \textit{WSDM 2012}, 673--682.\\[2pt]
Kempe, D., Kleinberg, J., \& Tardos, \'E. (2003). Maximizing the spread of influence through a social network. \textit{KDD '03}, 137--146.\\[2pt]
Lally, P., et al. (2010). How are habits formed: Modelling habit formation in the real world. \textit{European Journal of Social Psychology, 40}, 998--1009.\\[2pt]
Nowzari, C., Preciado, V. M., \& Pappas, G. J. (2016). Analysis and control of epidemics: A survey of spreading processes on complex networks. \textit{IEEE Control Systems Magazine, 36}, 26--46.\\[2pt]
Pastor-Satorras, R., \& Vespignani, A. (2001). Epidemic spreading in scale-free networks. \textit{Physical Review Letters, 86}, 3200--3203.\\[2pt]
Pastor-Satorras, R., Castellano, C., Van Mieghem, P., \& Vespignani, A. (2015). Epidemic processes in complex networks. \textit{Reviews of Modern Physics, 87}, 925--979.\\[2pt]
Preciado, V. M., et al. (2014). Optimal resource allocation for network protection against spreading processes. \textit{IEEE Transactions on Control of Network Systems, 1}, 99--108.\\[2pt]
Ryan, R. M., \& Deci, E. L. (2000). Intrinsic and extrinsic motivations: Classic definitions and new directions. \textit{Contemporary Educational Psychology, 25}, 54--67.\\[2pt]
Ugander, J., Backstrom, L., Marlow, C., \& Kleinberg, J. (2012). Structural diversity in social contagion. \textit{PNAS, 109}, 5962--5966.\\[2pt]
Watts, D. J. (2002). A simple model of global cascades on random networks. \textit{PNAS, 99}, 5766--5771.
\end{flushleft}

\end{document}
